% =============================================================================
% Conclusion
% =============================================================================
\section{Conclusion}\label{sec:conclusion}

We have investigated the practical identifiability of geometric and material
parameters of a fluid-filled viscoelastic oblate spheroidal shell from its
resonant frequencies.  The principal findings are:

\begin{enumerate}
  \item The equivalent-sphere model renders the inverse problem
    \emph{ill-posed}: the Jacobian of the frequency--parameter map is
    rank-deficient (condition number $\sim 10^{10}$), because all flexural
    modes scale identically with a single effective radius.
  \item The oblate spheroid (Rayleigh--Ritz) model yields a
    \emph{well-conditioned} inverse problem (condition number $\approx 83.5$),
    because different modes sample the anisotropic curvature differently.
  \item Newton--Raphson inversion recovers the semi-major axis~$a$,
    semi-minor axis~$c$, and Young's modulus~$E$ to sub-percent accuracy
    from three measured flexural frequencies, starting from a $10\%$-perturbed
    initial guess.
  \item Cram\'er--Rao lower bounds confirm that the minimum achievable
    parameter uncertainty scales linearly with measurement noise, with
    relative bounds of a few percent for $1\%$ frequency noise.
\end{enumerate}

In the language of Kac: one can hear the shape of an oblate organ, but not a
spherical one.  The practical implication is that aspherical forward models
are not merely more accurate --- they are \emph{essential} for solving the
inverse problem.
